\usepackage[margin=2.45cm,headheight=0pt,footskip=1.0cm]{geometry}
\usepackage{fontspec}
\setmainfont{texgyretermes-regular.otf}[BoldFont=texgyretermes-bold.otf,ItalicFont=texgyretermes-italic.otf,BoldItalicFont=texgyretermes-bolditalic.otf]
\setsansfont{texgyreheros-regular.otf}[BoldFont=texgyreheros-bold.otf,ItalicFont=texgyreheros-italic.otf]
\setmonofont{DejaVuSansMono.ttf}[Path=fonts/,Scale=0.88]
\IfFontExistsTF{SimSun}{\setCJKmainfont{SimSun}[AutoFakeBold=2.5]\setCJKmonofont{SimSun}}{\setCJKmainfont{FandolSong-Regular.otf}[BoldFont=FandolSong-Bold.otf]\setCJKmonofont{FandolSong-Regular.otf}}
\IfFontExistsTF{Heiti SC}{\setCJKsansfont{Heiti SC}}{\setCJKsansfont{FandolHei-Regular.otf}}
\usepackage{amsmath,amssymb,bm,mathtools}
\usepackage{graphicx,booktabs,array,tabularx,longtable,multirow}
\usepackage{subcaption,caption,float,placeins}
\usepackage{enumitem,xcolor,siunitx}
\usepackage{tikz}
\usetikzlibrary{arrows.meta,positioning,calc,fit}
\usepackage{fancyvrb,fvextra}
\usepackage{xurl}
\usepackage[unicode,hidelinks]{hyperref}
\usepackage{bookmark}
\hypersetup{pdftitle={基于多源数据与结构水热耦合的吉隆泥石流触发及传播研究},pdfauthor={},pdfsubject={赛道B论文初稿}}
\pagestyle{plain}
\linespread{1.25}
\setlength{\parindent}{2em}
\setlength{\parskip}{0.15em}
\setlength{\emergencystretch}{2em}
\setlength{\textfloatsep}{10pt plus 2pt minus 2pt}
\setlength{\floatsep}{8pt plus 2pt minus 2pt}
\setlength{\intextsep}{8pt plus 2pt minus 2pt}
\renewcommand{\topfraction}{0.9}
\renewcommand{\bottomfraction}{0.8}
\renewcommand{\textfraction}{0.08}
\renewcommand{\floatpagefraction}{0.75}
\setcounter{topnumber}{3}
\setcounter{totalnumber}{5}
\setcounter{tocdepth}{1}
\captionsetup{font=small,labelfont=bf,labelsep=quad}
\setlist{nosep,leftmargin=2em}
\sisetup{detect-all,group-minimum-digits=5}
\ctexset{section={format=\Large\bfseries,beforeskip=1.1ex plus .3ex,afterskip=.7ex},subsection={format=\large\bfseries,beforeskip=.9ex plus .2ex,afterskip=.5ex},subsubsection={format=\normalsize\bfseries}}
\newcommand{\dif}{\mathop{}\!\mathrm d}
\newcommand{\E}{\mathbb E}
\newcommand{\Prob}{\mathbb P}
\newcommand{\given}{\,\vert\,}
\newcommand{\AIcredit}{\par\begingroup\footnotesize 本节文字整理、程序与图表制作使用了OpenAI Codex辅助，贡献记录见文后声明及配套说明。\par\endgroup}
\newcolumntype{Y}{>{\raggedright\arraybackslash}X}
\newcommand{\sourcecode}[2]{\subsection*{#1}\begingroup\footnotesize\VerbatimInput[breaklines=true,breakanywhere=true,numbers=left,numbersep=4pt,fontsize=\scriptsize,tabsize=4]{#2}\endgroup}
